\documentclass[11pt,reqno]{amsart}

\usepackage[T1]{fontenc}
\usepackage[utf8]{inputenc}
\usepackage[english]{babel}
\usepackage{lmodern}
\usepackage{microtype}
\usepackage{amsmath,amssymb,amsthm,mathtools}
\usepackage{xcolor}
\usepackage{tikz}
\usepackage{hyperref}

\raggedbottom

\newtheorem{theorem}{Theorem}[section]
\newtheorem{proposition}[theorem]{Proposition}
\newtheorem{corollary}[theorem]{Corollary}
\newtheorem{lemma}[theorem]{Lemma}
\newtheorem{conjecture}[theorem]{Conjecture}
\newtheorem*{introtheorem}{Theorem}
\theoremstyle{definition}
\newtheorem{definition}[theorem]{Definition}
\newtheorem{example}[theorem]{Example}
\newtheorem{problem}[theorem]{Problem}
\newtheorem{remark}[theorem]{Remark}

\DeclareMathOperator{\code}{code}
\newcommand{\setN}{\mathbb{N}}
\newcommand{\setZ}{\mathbb{Z}}
\newcommand{\xvec}{\mathbf{x}}
\newcommand{\symS}{S}
\newcommand{\schubertS}{\mathfrak{S}}
\newcommand{\keypoly}{\kappa}
\newcommand{\stretchop}{\mathbin{\ast}}
\newcommand{\stretchperm}[2]{#1\stretchop #2}
\newcommand{\defin}[1]{%
  \relax\ifmmode%
  \textcolor{blue}{#1}%
  \else\textcolor{blue}{\emph{#1}}%
  \fi%
}
\newcommand{\diagramcell}[3][gray!25]{%
  \draw[thick,fill=#1] ({#2-1},{#3-1}) rectangle (#2,#3);%
}
\newcommand{\cellDiagram}[4][.42]{%
\begin{tikzpicture}[scale=#1,baseline=(current bounding box.south)]
  \draw[thin,gray!35] (0,0) grid (#2,#3);
  #4
  \draw[very thick] (0,0) rectangle (#2,#3);
\end{tikzpicture}%
}
\newcommand{\rotheDiagram}[2][]{%
\begin{tikzpicture}[baseline=(current bounding box.south),#1]
  \def\perm{#2}
  \def\rotheN{0}
  \foreach \val [count=\i] in \perm {\xdef\rotheN{\i}}
  \draw[thin,gray!30] (0,0) grid (\rotheN,\rotheN);
  \foreach \wi [count=\i] in \perm {
    \foreach \wj [count=\j] in \perm {
      \ifnum\i<\j
        \ifnum\wi>\wj
          \draw[thick,fill=gray!25,fill opacity=.72]
            (\wj-1,\i-1) rectangle (\wj,\i);
        \fi
      \fi
    }
  }
  \foreach \wi [count=\i] in \perm {
    \draw[line width=1.2pt]
      (\wi-.5,\i-.5) -- (\rotheN,\i-.5);
    \draw[line width=1.2pt]
      (\wi-.5,\i-.5) -- (\wi-.5,\rotheN);
    \fill (\wi-.5,\i-.5) circle (.15);
  }
  \draw[very thick] (0,0) rectangle (\rotheN,\rotheN);
\end{tikzpicture}%
}
\newcommand{\rotheDiagramWithBlocks}[2][]{%
\begin{tikzpicture}[baseline=(current bounding box.south),#1]
  \def\perm{#2}
  \def\rotheN{0}
  \foreach \val [count=\i] in \perm {\xdef\rotheN{\i}}
  \draw[thin,gray!30] (0,0) grid (\rotheN,\rotheN);
  \foreach \wi [count=\i] in \perm {
    \foreach \wj [count=\j] in \perm {
      \ifnum\i<\j
        \ifnum\wi>\wj
          \def\cellfill{gray!25}
          \ifnum\wj>1
            \ifnum\wj<5
              \def\cellfill{blue!18}
            \fi
          \fi
          \ifnum\wj>6
            \ifnum\wj<9
              \def\cellfill{green!18}
            \fi
          \fi
          \draw[thick,fill=\cellfill,fill opacity=.82]
            (\wj-1,\i-1) rectangle (\wj,\i);
        \fi
      \fi
    }
  }
  \foreach \wi [count=\i] in \perm {
    \draw[line width=1.2pt]
      (\wi-.5,\i-.5) -- (\rotheN,\i-.5);
    \draw[line width=1.2pt]
      (\wi-.5,\i-.5) -- (\wi-.5,\rotheN);
    \fill (\wi-.5,\i-.5) circle (.15);
  }
  \draw[very thick] (0,0) rectangle (\rotheN,\rotheN);
\end{tikzpicture}%
}

\title[Stretched Schubert and Key Coefficients]
{Polynomiality of Stretched Schubert Structure Constants and Key Coefficients}
\author{Per Alexandersson}
\address{Department of Mathematics, Stockholm University,
  SE-106 91 Stockholm, Sweden}
\email{per.w.alexandersson@gmail.com}
\subjclass[2020]{Primary 05E05; Secondary 14N15, 05A15}
\keywords{Schubert polynomial, key polynomial, Schubert structure constant,
  Demazure operator, vector partition function, Rothe diagram}
\date{}

\begin{document}

\begin{abstract}
We prove that monomial coefficients in affine families of key and Schubert
polynomials are eventually polynomial.  The proof combines Demazure operators
with vector partition functions and, in the Schubert case, P.~Magyar's
orthodontic formula.  These coefficient results extend to finite products.  Using
M.~Watanabe's Schubert duality, we deduce that stretched Schubert structure
constants are eventually polynomial, proving a conjecture of I.~Pak and
Z.~Slonim.  For key polynomials, this resolves the polynomiality part of a
conjecture of P.~Alexandersson and E.~Alhajjar.
\end{abstract}

\maketitle

\section{Introduction}
\label{sec:introduction}

Schubert polynomials form a distinguished basis of the polynomial ring, and
their multiplication is governed by the \defin{Schubert structure constants}.
For permutations \(u,v\), we define
\(\defin{c_{u,v}^w}\) by
\[
\schubertS_u\schubertS_v
=\sum_w c_{u,v}^w\schubertS_w.
\]
Let \(\defin{\code(u)}\) denote the Lehmer code of \(u\), and let
\(\defin{\ell(u)}\) denote its inversion number.  For a positive integer
\(N\), we define
\[
\defin{\stretchperm{N}{u}}\coloneqq\code^{-1}\bigl(N\code(u)\bigr).
\]
If \(N\code(u)=(c_1,\dotsc,c_h)\), we append \(t\) trailing zeroes with
\(c_i\leq h+t-i\) for all \(i\); the resulting finite sequence is then a
Lehmer code.  Increasing \(t\) only appends fixed points to the corresponding
permutation.
I.~Pak and Z.~Slonim proved that the resulting stretched Schubert structure
constants are eventually quasipolynomial~\cite{PakSlonim2026x}, and
conjectured that the period is always one.  We prove their conjecture by first
considering affine monomial coefficients of stretched Schubert polynomials.

For
\[
\defin{\xvec}\coloneqq(x_1,\dotsc,x_h),
\qquad
\defin{\xvec^a}\coloneqq x_1^{a_1}\dotsm x_h^{a_h}
\quad\text{for }a\in\setZ^h,
\]
we write \(\defin{[\xvec^a]f}\) for the coefficient of the monomial
\(\xvec^a\) in \(f\).
Classical Kostka numbers are Schur monomial coefficients.  Following the
Schubert--Kostka matrix terminology of A.~Postnikov and R.~Stanley
\cite[Sec.~17]{PostnikovStanley2008}, we call the analogous monomial
coefficients of Schubert polynomials Schubert--Kostka coefficients; similarly,
we call the monomial coefficients of key polynomials key--Kostka coefficients.
We write
\(\defin{\setN}\coloneqq\{0,1,2,\dotsc\}\).  A function of \(N\) is
\defin{eventually polynomial} if there is a polynomial \(P(N)\) such that the
function agrees with \(P(N)\) for all \(N\gg0\).

The remaining definitions used in these statements are recalled below.  The
same monomial-coefficient question can be asked for key polynomials, or
type~\(A\) Demazure characters.

\begin{introtheorem}[{\hyperref[cor:pureKeyKostka]
  {see Corollary~\ref*{cor:pureKeyKostka}}}]
  Let \(\alpha,\nu\) be weak compositions of the same total degree.
  Then the map
\[
N\longmapsto
[\xvec^{N\nu}]\keypoly_{N\alpha}(\xvec)
\]
is a polynomial in \(N\).
\end{introtheorem}

The key--Kostka theorem settles the polynomiality part of a
conjecture from earlier work with E.~Alhajjar
\cite[Conj.~5.1]{AlexanderssonAlhajjar2018}.  The conjectured nonnegativity of
the resulting polynomial remains open.  Such positivity is unknown even for
classical stretched Littlewood--Richardson and Kostka polynomials
\cite{KingTolluToumazet2004}.  Another recent approach to stretched Kostka
coefficients, using partition-division maps on symmetric functions, is given
in~\cite{AlexanderssonDai2026x}.

The Schubert analogue is eventually polynomial along every affine weight ray.

\begin{introtheorem}[{\hyperref[thm:affineSchubertKostka]
  {see Theorem~\ref*{thm:affineSchubertKostka}}}]
Fix a permutation \(u\) and an integer \(h\) such that
\(\code(u)\in\setN^h\).  Let \(\nu,\eta\in\setZ^h\), and suppose that
\(N\nu+\eta\) is a weak composition and has total degree \(N\ell(u)\) for all
sufficiently large \(N\).  Then
\[
N\longmapsto
[\xvec^{N\nu+\eta}]\schubertS_{\stretchperm{N}{u}}(\xvec)
\]
is eventually polynomial.
\end{introtheorem}

Schubert duality then gives the main theorem.

\begin{introtheorem}[{\hyperref[cor:stretchedSchubertStructureConstants]
  {see Corollary~\ref*{cor:stretchedSchubertStructureConstants}}}]
For fixed permutations \(u,v,w\), the stretched Schubert structure constant
\[
N\longmapsto
c_{\stretchperm{N}{u},\stretchperm{N}{v}}^{\stretchperm{N}{w}}
\]
is eventually polynomial.
\end{introtheorem}

The two coefficient theorems reduce coefficient extraction from fixed Demazure
expressions to type~\(A\) vector partition functions.  For Schubert
polynomials, P.~Magyar's orthodontic formula supplies such an expression once
the affine column multiplicities have stabilized
\cite[Prop.~15]{Magyar1998SchubertBottSamelson}.  M.~Watanabe's duality formula
then expresses each Schubert structure constant as a fixed alternating sum of
monomial coefficients~\cite[Lem.~7.4]{Watanabe2016}.

The paper is organized as follows.  Section~\ref{sec:dividedDifferenceLemma}
gives the vector-partition-function lemma used throughout.  We apply it to key
polynomials in Section~\ref{sec:keyKostka}.  Sections
\ref{sec:stretchingRotheDiagrams} and~\ref{sec:orthodontia} establish the
affine orthodontic description of stretched Rothe diagrams and prove affine
Schubert--Kostka polynomiality.  Finally, Section~\ref{sec:products} treats
products and Schubert structure constants.

\section{A divided-difference lemma}
\label{sec:dividedDifferenceLemma}

We write
\(\defin{s_i}\) for the transposition of \(x_i\) and \(x_{i+1}\).  Let
\[
\defin{\pi_i(f)}
\coloneqq
\frac{x_i f-x_{i+1}s_i(f)}{x_i-x_{i+1}}.
\]
The operator \(\pi_i\) is the \defin{isobaric divided difference}, or
\defin{Demazure operator}.

The next lemma is the common polynomiality mechanism for keys and Schubert
polynomials.

\begin{lemma}[Fixed Demazure expressions]
\label{lem:fixedDemazureExpressions}
Fix \(h\).  Suppose that, for all sufficiently large \(N\), the Laurent
polynomial \(F_N\) is obtained from \(1\) by a fixed finite sequence of the
following operations:
\begin{enumerate}
\item multiply by a Laurent monomial
  \(\xvec^{Na+b}\), where \(a,b\in\setZ^h\) are fixed;
\item apply a fixed Demazure operator \(\pi_i\).
\end{enumerate}
Then, for fixed \(\nu,\eta\in\setZ^h\), the map
\[
N\longmapsto
[\xvec^{N\nu+\eta}]F_N
\]
is eventually polynomial in \(N\).  The same conclusion holds for a product
of finitely many such expressions.
\end{lemma}

\begin{proof}
Let \(\defin{\alpha_i}\) be the \(i\)-th simple root of type~\(A\).  Thus
\(\alpha_i\) has a \(1\) in position \(i\), a \(-1\) in position \(i+1\),
and zeroes elsewhere; it is the difference of two consecutive unit vectors.
On Laurent series we have
\[
\pi_i(f)
=
\frac{f}{1-\xvec^{-\alpha_i}}
-
\frac{\xvec^{-\alpha_i}s_i(f)}{1-\xvec^{-\alpha_i}}.
\]
Starting with a Laurent monomial and iterating this identity gives a finite
sum of terms
\begin{equation}
\label{eq:rootDenominatorExpansion}
C\xvec^{N\beta+\gamma}
\prod_{\rho\in\Phi}
\frac{1}{1-\xvec^{-\rho}},
\end{equation}
where \(C\in\setZ\), the vectors \(\beta,\gamma\in\setZ^h\) are fixed, and
\(\Phi\) is a finite \emph{multiset} of type~\(A\) roots.  If a denominator
contains a negative root, we use
\[
\frac{1}{1-\xvec^\rho}
=-
\frac{\xvec^{-\rho}}{1-\xvec^{-\rho}}
\]
to make it positive.  Thus all denominator roots in
\eqref{eq:rootDenominatorExpansion} may be taken positive.

The coefficient of \(\xvec^{N\nu+\eta}\) in one such term is \(C\) times
\[
\#\left\{(m_\rho)_{\rho\in\Phi}\in\setN^\Phi:
  \sum_{\rho\in\Phi}m_\rho\rho
  =N(\beta-\nu)+(\gamma-\eta)\right\}.
\]
This is a vector partition function.
A matrix whose columns are positive type~\(A\) roots
is a submatrix of a directed incidence matrix and is therefore totally
unimodular.  Repeated columns preserve total unimodularity.  By Sturmfels'
chamber theorem~\cite{Sturmfels1995}, the corresponding vector partition
function is polynomial on every relatively open face of its chamber complex.
The chamber complex is finite, and the signs of its finitely many defining
linear forms on the affine ray \(N(\beta-\nu)+(\gamma-\eta)\) are eventually
constant.  Hence the ray eventually remains in one face, so each term, and
therefore their finite sum, is eventually polynomial.

Multiplying finitely many expressions merely concatenates the multisets of
roots in the denominators, so the same argument proves the product statement.
\end{proof}

\section{Key-Kostka coefficients}
\label{sec:keyKostka}

Key polynomials, introduced by A.~Lascoux and M.-P.~Sch\"utzenberger
\cite{Lascoux1990Keys}, are the type~\(A\) Demazure characters.  Combinatorial
models include S.~K.~Mason's semi-skyline augmented fillings, which connect
nonsymmetric Macdonald polynomials with Demazure atoms and operators
\cite{Mason2009}, and the Kohnert tableaux of S.~H.~Assaf and D.~Searles
\cite{AssafSearles2018,AssafSearles2019}.

A \defin{weak composition} is a finite sequence of nonnegative integers; we
append trailing zeroes whenever necessary.  We use the divided-difference
definition of the key polynomial \(\keypoly_\alpha\).  If \(\alpha\) is weakly
decreasing, then
\[
\defin{\keypoly_\alpha}\coloneqq\xvec^\alpha.
\]
If \(\alpha_i<\alpha_{i+1}\), then
\[
\defin{\keypoly_\alpha}\coloneqq\pi_i\keypoly_{s_i\alpha}.
\]
This recursion is well-founded: each such step decreases by one the number of
pairs \(p<q\) with \(\alpha_p<\alpha_q\).
For fixed \(h\), the family
\[
\{\keypoly_\alpha:\alpha\in\setN^h\}
\]
is a basis of \(\mathbb Q[x_1,\dotsc,x_h]\).
For a permutation \(w\) with a reduced decomposition
\(w=s_{i_1}\dotsm s_{i_\ell}\), define
\[
\defin{\pi_w}
\coloneqq
\pi_{i_1}\dotsm\pi_{i_\ell}.
\]
The Demazure relations imply that \(\pi_w\) is independent of the chosen
reduced decomposition.
Equivalently, if \(\lambda\) is the decreasing rearrangement of \(\alpha\)
and \(w\) is the shortest permutation satisfying \(w\lambda=\alpha\), then
\(\keypoly_\alpha=\pi_w\xvec^\lambda\).

\begin{theorem}[Affine key-Kostka polynomiality]
\label{thm:affineKeyKostka}
Fix \(\alpha,\nu\in\setN^h\) and \(\beta,\eta\in\setZ^h\).  Suppose that
\(N\alpha+\beta\) and
\(N\nu+\eta\) are weak compositions of the same total degree for all
sufficiently large \(N\).  Then
\[
N\longmapsto
[\xvec^{N\nu+\eta}]\keypoly_{N\alpha+\beta}(\xvec)
\]
is eventually polynomial.
\end{theorem}

\begin{proof}
For \(N\gg0\), the relative order of the parts of \(N\alpha+\beta\) is fixed.
Its decreasing rearrangement therefore has the form \(N\lambda+\lambda_0\),
where \(\lambda,\lambda_0\) are fixed, and the shortest permutation \(w\) satisfying
\[
w(N\lambda+\lambda_0)=N\alpha+\beta
\]
is fixed.  Choose a reduced word for \(w\).  The expression
\[
\keypoly_{N\alpha+\beta}
=\pi_w\xvec^{N\lambda+\lambda_0}
\]
now satisfies Lemma~\ref{lem:fixedDemazureExpressions}, which proves the
claim.
\end{proof}

\begin{corollary}
\label{cor:pureKeyKostka}
For weak compositions \(\alpha,\nu\) of the same total degree,
\[
N\longmapsto
[\xvec^{N\nu}]\keypoly_{N\alpha}(\xvec),
\qquad N\geq0,
\]
is a polynomial in \(N\).
\end{corollary}

\begin{proof}
Theorem~\ref{thm:affineKeyKostka}, with \(\beta=\eta=0\), gives a polynomial
for \(N\gg0\).  The Kogan-face model for Demazure characters realizes the
same function as an Ehrhart quasipolynomial
\cite{Kogan2005,KiritchenkoSmirnovTimorin2010}; see also
\cite{AlexanderssonAlhajjar2018}.  A quasipolynomial that agrees with one
polynomial for all sufficiently large \(N\) agrees with that polynomial for
every \(N\geq0\).
\end{proof}

\begin{proposition}[Products of affine keys]
\label{prop:productsAffineKeys}
Fix \(h\), and let \(\alpha_j\in\setN^h\) and
\(\beta_j\in\setZ^h\), for \(1\leq j\leq s\), such that
\(N\alpha_j+\beta_j\) is a weak composition for \(N\gg0\).  If
\(\nu\in\setN^h\), \(\eta\in\setZ^h\), and \(N\nu+\eta\) is eventually
nonnegative, then
\[
[\xvec^{N\nu+\eta}]
\prod_{j=1}^s\keypoly_{N\alpha_j+\beta_j}(\xvec)
\]
is eventually polynomial.
\end{proposition}

\begin{proof}
Use the fixed reduced word from the proof of
Theorem~\ref{thm:affineKeyKostka} for each factor, and apply the product part
of Lemma~\ref{lem:fixedDemazureExpressions}.  The claim follows.
\end{proof}

Theorem~\ref{thm:affineKeyKostka} extends the polynomiality result from joint
work with E.~Kantarc{\i} O{\u{g}}uz from stretched row-flagged skew Schur
polynomials and cylindric Schur functions to affine families of key
polynomials~\cite{AlexanderssonOguz2023x}.

\section{Stretching Rothe diagrams}
\label{sec:stretchingRotheDiagrams}

Let \(\defin{\symS_m}\) denote the symmetric group on \(m\) letters.
For \(w\in\symS_m\), its \defin{Rothe diagram} is
\[
\defin{D(w)}
\coloneqq
\{(i,j):1\leq i,j\leq m,\ j<w(i),\ i<w^{-1}(j)\}.
\]
We use matrix coordinates: \(i\) is the row index and \(j\) is the column
index.  Thus the cells in row \(i\) are the pairs \((i,j)\in D(w)\).  The row
lengths form its \defin{Lehmer code},
\[
\defin{\code(w)_i}
\coloneqq
\#\{j>i:w(j)<w(i)\}.
\]
Conversely, if \(\gamma=(\gamma_1,\dotsc,\gamma_h)\) is a weak composition
and \(t\) satisfies \(\gamma_i\leq h+t-i\) for all \(i\), then
\((\gamma_1,\dotsc,\gamma_h,0^t)\) is the Lehmer code of a unique
permutation in \(\symS_{h+t}\).

We regard \(D(w)\) as a word of columns.  The \defin{type of a column} is the
set of rows occupied in that column.

\begin{example}[The affine column blocks for \(321654\)]
\label{ex:rothe321654}
Let \(\defin{u}\coloneqq321654\), so that
\(\code(u)=(2,1,0,2,1,0)\).  For \(N\geq3\),
the nonempty columns of \(D(\stretchperm{N}{u})\), read from left to right,
have types
\begin{equation}
\label{eq:321654ColumnWord}
\{1,2\},\quad
\{1,2,4,5\}^{N-1},\quad
\{1\},\quad
\{1,4,5\},\quad
\{1,4\}^{N-2},\quad
\{4\}^{2}.
\end{equation}
For \(N=4\), we have
\[
\stretchperm{4}{u}=(9,5,1,12,7,2,3,4,6,8,10,11),
\]
and its Rothe diagram is shown below.
\[
\rotheDiagramWithBlocks[scale=.3125]{9,5,1,12,7,2,3,4,6,8,10,11}
\]
The three blue columns have type \(\{1,2,4,5\}\), while the two green
columns have type \(\{1,4\}\).  These are the two column blocks whose
lengths grow with \(N\).
\end{example}

A permutation is \defin{\(k\)-Grassmannian} if it has no descents except
possibly at position \(k\).  For Grassmannian permutations, code stretching
has the usual interpretation as stretching a Ferrers shape.

\begin{lemma}[Grassmannian stretching]
\label{lem:grassmannianStretching}
Let \(\lambda=(\lambda_1,\dotsc,\lambda_k)\) be a partition, and let \(u\) be
the permutation with Lehmer code \((\lambda_k,\dotsc,\lambda_1)\).  Then
\(u\) is \(k\)-Grassmannian and
\(\schubertS_u(x_1,\dotsc,x_k)=\defin{s_\lambda(x_1,\dotsc,x_k)}\), the Schur
polynomial indexed by \(\lambda\).  Moreover,
\(\stretchperm{N}{u}\) is the Grassmannian permutation associated with the
stretched shape \(N\lambda\), and
\[
\schubertS_{\stretchperm{N}{u}}(x_1,\dotsc,x_k)
=
s_{N\lambda}(x_1,\dotsc,x_k).
\]
\end{lemma}

\begin{proof}
The Schubert--Schur identity for Grassmannian permutations is standard; see,
for example,~\cite{ReinerShimozono1995}.  The remaining assertion follows from
\[
\code(\stretchperm{N}{u})
=
(N\lambda_k,\dotsc,N\lambda_1).
\]
\end{proof}

Pak and Slonim prove a stability statement for pipe-dream polytopes under code
scaling~\cite[Rem.~4.4]{PakSlonim2026x}.  The next proposition gives the
corresponding finite-block description for Rothe diagram columns, with an
additional affine shift of the code.

\begin{proposition}[Affine column multiplicities]
\label{prop:affineColumnMultiplicities}
Fix \(\alpha\in\setN^h\) and \(\beta\in\setZ^h\), and suppose that
\(N\alpha+\beta\) is nonnegative for \(N\gg0\).  Let \(w(N)\) be the
permutation with Lehmer code \(N\alpha+\beta\), after appending sufficiently
many trailing zeroes.  For \(N\gg0\), every nonempty column of \(D(w(N))\)
has its rows in \(\{1,\dotsc,h\}\), and each column type
\(C\subseteq\{1,\dotsc,h\}\) occurs with multiplicity
\[
a_C N+b_C,
\]
where \(a_C,b_C\in\setZ\) are independent of \(N\).
\end{proposition}

\begin{proof}
Write \(\alpha=(a_1,\dotsc,a_h)\) and
\(\beta=(b_1,\dotsc,b_h)\).  Since the code entries after position \(h\)
are zero, every row of \(D(w(N))\) with index greater than \(h\) is empty.
In the inverse Lehmer-code construction, \(w(N)(i)\) is the
\((Na_i+b_i+1)\)-st positive integer not used in the first \(i-1\) positions.
Thus, for some \(r\in\{0,\dotsc,i-1\}\),
\[
w(N)(i)=Na_i+b_i+1+r.
\]
Induction on \(i\) shows that this \(r\) is constant for \(N\gg0\).  Indeed,
assume that the earlier values \(w(N)(1),\dotsc,w(N)(i-1)\) are affine in
\(N\).  For a fixed \(r\), the value \(Na_i+b_i+1+r\) is selected precisely
when exactly \(r\) earlier values are smaller, and no earlier value is equal
to it.  These are finitely many comparisons between affine functions.  Since
two affine functions are either identical or meet at at most one value of
\(N\), each comparison has an eventually fixed truth value.  Since exactly
one \(r\) is selected for each \(N\), the selected \(r\) is eventually
constant.  Hence the first \(h\) values of \(w(N)\) are affine in \(N\), and
their strict relative order is fixed.

If a column index \(j\) is not one of these values, its column type is
\[
\{i\in\{1,\dotsc,h\}:j<w(N)(i)\}.
\]
Indeed, if \(j<w(N)(i)\) and \(j\) has not occurred among the first \(h\)
values, then \(j\) is unused before position \(i\), so \(i<w(N)^{-1}(j)\).
The type is empty when \(j\) exceeds all first \(h\) values, and it is
constant as \(j\) runs between consecutive values of
\(w(N)(1),\dotsc,w(N)(h)\).  Since the strict relative order of these affine
values is fixed for large \(N\), each gap is either eventually empty or has
size equal to the difference of its endpoints minus one.  In either case its
size is affine in \(N\).  The initial gap has size
\(\min_i w(N)(i)-1\), also affine.

It remains to account for the columns \(j=w(N)(s)\) with \(s\leq h\).  Such a
column has type
\[
\{i<s:w(N)(s)<w(N)(i)\},
\]
which is fixed by the relative order.  Summing the gap sizes and these
individual columns proves the claim.
\end{proof}

\section{Orthodontia and Schubert-Kostka coefficients}
\label{sec:orthodontia}

Postnikov and Stanley define the \defin{Schubert--Kostka matrix} as the matrix
of monomial coefficients \(K_{u,a}=[\xvec^a]\schubertS_u\)
\cite[Sec.~17]{PostnikovStanley2008}.  Pak and Slonim use the term
\defin{Schubert--Kostka number} for these coefficients
\cite[Sec.~2.6]{PakSlonim2026x}.

\subsection{Magyar's orthodontic formula}

We recall Magyar's orthodontic construction and formula
\cite[Sec.~4.3 and Prop.~15]{Magyar1998SchubertBottSamelson}.

For \(r\geq1\), put
\[
\defin{[r]}\coloneqq\{1,2,\dotsc,r\},
\qquad
\defin{\xvec_r}\coloneqq x_1x_2\dotsm x_r.
\]
In this construction, a Rothe diagram is regarded as a multiset of nonempty
column types.  For each \(r\), remove all columns of type \([r]\), and let
\(\defin{k_r}\) be their number.  If any columns remain, order their types
lexicographically.  In the first remaining type, choose its
smallest \defin{missing tooth} \(i\): row \(i\) is absent and row \(i+1\) is
present.
Interchange rows \(i\) and \(i+1\) in every remaining column.  Remove the
columns that have become \([i]\), record their number as \(m\), and
repeat.  This gives the \defin{orthodontic word}
\[
\defin{\mathbf{i}}\coloneqq(i_1,\dotsc,i_\ell)
\]
and multiplicities \(\defin{m_1,\dotsc,m_\ell}\).  Each iteration contributes
one entry to \(\mathbf{i}\) and performs one row interchange; thus \(\ell\) is
the number of row interchanges.

Magyar proves that the process terminates for every fixed diagram, and his
formula states that
\begin{equation}
\label{eq:magyarOrthodontia}
\schubertS_w
=
\left(\prod_{r=1}^h\xvec_r^{k_r}\right)
\bigl(
\pi_{i_1}\circ \xvec_{i_1}^{m_1}
\circ\pi_{i_2}\circ \xvec_{i_2}^{m_2}
\circ\dotsm
\circ\pi_{i_\ell}\circ \xvec_{i_\ell}^{m_\ell}
\bigr)(1).
\end{equation}
The composition is read from right to left, with monomials acting by
multiplication.

Magyar proves that \eqref{eq:magyarOrthodontia} can be written as a sum of
\(2^\ell\) rational functions
\cite[Thm.~3]{Magyar1998SchubertBottSamelson}.
After rewriting each negative denominator root as a positive one, every
summand has the form
\begin{equation}
\label{eq:magyarRationalTerm}
\epsilon\,\xvec^\gamma
\prod_{\rho\in\Phi}\frac{1}{1-\xvec^{-\rho}},
\end{equation}
where \(\epsilon\in\{1,-1\}\), \(\gamma\in\setZ^h\), and \(\Phi\) is a
multiset of positive type~\(A\) roots.  This is the root-denominator form
used in the proof of Lemma~\ref{lem:fixedDemazureExpressions}.

\subsection{Affine orthodontic data}

To apply Lemma~\ref{lem:fixedDemazureExpressions} uniformly in \(N\), it
remains to show that the operator word is fixed and that the exponents are
affine in \(N\).

\begin{lemma}[Stable orthodontic data]
\label{lem:stableOrthodonticData}
Let \(w(N)\) be a permutation family whose Lehmer code is affine in \(N\),
as in Proposition~\ref{prop:affineColumnMultiplicities}.  For \(N\gg0\),
the orthodontic word \(\mathbf{i}\) for \(D(w(N))\) is independent of \(N\),
and all multiplicities \(k_r(N)\) and \(m_j(N)\) are affine in \(N\).
\end{lemma}

\begin{proof}
By Proposition~\ref{prop:affineColumnMultiplicities}, the initial
column-type multiplicities are affine in \(N\).  If an integer-valued affine
function is nonnegative for all large \(N\), then it is either identically
zero, eventually a positive constant, or eventually tends to \(+\infty\).
Thus the initial support is fixed for large \(N\).  The multiplicities
\(k_r(N)\) of the removed columns \([r]\) are therefore affine.

We now run Magyar's algorithm.  Suppose at some stage that the support of the
remaining column types is fixed and that all their multiplicities are affine.
The lexicographic order and smallest-tooth rule choose the same next index
\(i\) for all sufficiently large \(N\).  Swapping rows \(i\) and \(i+1\)
sends each fixed column type to a fixed column type, so the new
multiplicities are sums of affine functions.  The recorded multiplicity
\(m(N)\) of the columns that have become \([i]\) is affine, and the remaining
support is again eventually fixed by the same positivity argument.

Fix one sufficiently large \(N_0\).  Inducting through the finitely many steps
of Magyar's algorithm for \(N_0\), the preceding paragraph shows that, after
possibly increasing the lower bound on \(N\), the same support sequence and
the same choices occur for every \(N\).  The support is empty at the final
stage for \(N_0\), hence also for all large \(N\).
Thus the successive supports, the orthodontic word, and the number of steps
are eventually independent of \(N\), while all recorded multiplicities are
affine in \(N\).
\end{proof}

\begin{example}[Orthodontia for \(321654\)]
\label{ex:orthodontia321654}
Continue with \(u=321654\) and the column word in
\eqref{eq:321654ColumnWord}.  For \(N\geq3\), orthodontia yields
\[
\mathbf{i}=(3,4,2,3,1),
\qquad
(m_1,\dotsc,m_5)=(0,N-1,N-2,1,2).
\]
Consequently,
\[
\schubertS_{\stretchperm{N}{321654}}
=
\xvec_1\xvec_2
\bigl(
\pi_3\circ\pi_4\circ \xvec_4^{N-1}
\circ\pi_2\circ \xvec_2^{N-2}
\circ\pi_3\circ \xvec_3
\circ\pi_1\circ \xvec_1^2
\bigr)(1).
\]
The operator word has length five and is independent of \(N\).  The two long
column blocks appear only in the powers \(N-1\) and \(N-2\).
\end{example}

\begin{theorem}[Affine Schubert-Kostka polynomiality]
\label{thm:affineSchubertKostka}
Fix \(\alpha,\nu\in\setN^h\) and \(\beta,\eta\in\setZ^h\), and let \(w(N)\)
be the permutation whose Lehmer code is obtained from \(N\alpha+\beta\) by
appending enough trailing zeroes to satisfy the Lehmer-code inequalities.
Suppose that \(N\alpha+\beta\) and
\(N\nu+\eta\) are nonnegative and have the same total degree for all
sufficiently large \(N\).  Then
\[
N\longmapsto
[\xvec^{N\nu+\eta}]\schubertS_{w(N)}(\xvec)
\]
is eventually polynomial.
\end{theorem}

\begin{proof}
By Lemma~\ref{lem:stableOrthodonticData}, the word in
\eqref{eq:magyarOrthodontia} is fixed for \(N\gg0\), and all exponents
\(k_r(N)\) and \(m_j(N)\) are affine in \(N\).  Hence
\eqref{eq:magyarOrthodontia} is a fixed Demazure expression of the form
covered by Lemma~\ref{lem:fixedDemazureExpressions}, which proves the claim.
\end{proof}

\subsection{Stable key expansions}

The preceding argument proves coefficient polynomiality.  Schubert polynomials
are key positive; this follows, for example, from the flagged
Littlewood--Richardson rule of Reiner and Shimozono~\cite{ReinerShimozono1995}
and from Assaf's Kohnert model~\cite{Assaf2021}.  The following example
suggests that this positive key expansion has a stronger stability property
under stretching.

\begin{example}[A stable key expansion]
\label{ex:stableKeyExpansion321654}
For the permutation \(u=321654\) from Example~\ref{ex:orthodontia321654},
the key expansion stabilizes at \(N=3\).  Omitting trailing zeroes, for
\(N\geq3\) we have
\[
\begin{aligned}
\schubertS_{\stretchperm{N}{321654}}
={}&
\keypoly_{(2N,N,0,2N,N)}
+\keypoly_{(2N,N+1,0,2N,N-1)}\\
&+\keypoly_{(2N,2N+1,0,N-1,N)}
+\keypoly_{(2N+1,N,0,2N-1,N)}\\
&+\keypoly_{(2N+1,N+1,0,2N-1,N-1)}
+\keypoly_{(2N+2,N,0,2N-2,N)}\\
&+\keypoly_{(2N+2,N+1,0,2N-2,N-1)}
+\keypoly_{(2N+2,2N-1,0,N-1,N)}.
\end{aligned}
\]
Thus the expansion has eight terms with constant coefficient \(1\), and every
key index is affine in \(N\).
\end{example}

\begin{conjecture}[Stable affine key expansion]
\label{conj:stableAffineKeyExpansion}
For every permutation \(u\), there exist positive integers \(c_1,\dotsc,c_t\)
and affine weak compositions
\[
N\gamma_j+\delta_j,
\qquad 1\leq j\leq t,
\]
such that, for all sufficiently large \(N\),
\[
\schubertS_{\stretchperm{N}{u}}
=
\sum_{j=1}^t
c_j\keypoly_{N\gamma_j+\delta_j}.
\]
In particular, both the number of terms and their coefficients eventually
stabilize, while the key indices are affine in \(N\).
\end{conjecture}

\begin{remark}
Conjecture~\ref{conj:stableAffineKeyExpansion} is stronger than
Theorem~\ref{thm:affineSchubertKostka}.  Magyar's formula fixes the number of
rational summands, but it does not fix the number of Demazure components in
the positive key expansion.  Example~\ref{ex:stableKeyExpansion321654}
exhibits the behavior asserted in the conjecture, but the divided-difference
argument does not imply this stronger conclusion.
\end{remark}

For a partition \(\lambda\) and a weakly increasing row flag
\(\mathbf b=(b_1,\dotsc,b_{\ell(\lambda)})\) with \(b_i\geq i\), write
\(\defin{s_\lambda(\mathbf b)}\) for the corresponding flagged Schur
polynomial.

\begin{conjecture}[Flagged Schur analogue]
\label{conj:stableFlaggedSchurKeyExpansion}
Let \(N\lambda+\mu\) be an affine family of partitions with a fixed number of
rows, and fix such a row flag \(\mathbf b\).  Then the key expansion of
the flagged Schur polynomial \(s_{N\lambda+\mu}(\mathbf b)\) eventually has a
fixed number of terms and fixed coefficients, with every key index affine in
\(N\).
\end{conjecture}

\begin{remark}[Flagged Schur functions]
For vexillary permutations, Schubert polynomials are flagged Schur functions,
and M.~Wachs' symmetrizing-operator formula puts affine families with a fixed
flag in the form covered by Lemma~\ref{lem:fixedDemazureExpressions}
\cite{Wachs1985}.  Their affine monomial coefficients are therefore eventually
polynomial.  This coefficientwise conclusion does not imply
Conjecture~\ref{conj:stableFlaggedSchurKeyExpansion}.
\end{remark}

M.~Shimozono and T.~Yu proved that Grothendieck polynomials expand positively
in Lascoux polynomials~\cite{ShimozonoYu2021}, which motivates the following
\(K\)-theoretic analogue.  We retain the connective-\(K\) parameter \(\beta\),
so that \(\beta=0\) recovers Schubert and key polynomials, and write
\(\defin{\mathfrak G_w^{(\beta)}}\) and
\(\defin{\mathfrak L_\alpha^{(\beta)}}\) for the corresponding
connective-\(K\) Grothendieck and Lascoux polynomials.

\begin{conjecture}[Stable Grothendieck--Lascoux expansion]
\label{conj:stableGrothendieckLascouxExpansion}
For every permutation \(u\), there exist positive integers \(c_1,\dotsc,c_t\),
nonnegative integers \(e_1,\dotsc,e_t\), and affine weak compositions
\(N\gamma_j+\delta_j\), for \(1\leq j\leq t\), such that, for \(N\gg0\),
\[
\mathfrak G_{\stretchperm{N}{u}}^{(\beta)}
=
\sum_{j=1}^t
c_j\beta^{e_j}\mathfrak L_{N\gamma_j+\delta_j}^{(\beta)}.
\]
\end{conjecture}

\section{Products and Schubert structure constants}
\label{sec:products}

The divided-difference argument is stable under finite products.  We record
the affine form needed for the subsequent extraction of Schubert structure
constants.

\begin{proposition}[Affine product Schubert-Kostka polynomiality]
\label{prop:productSchubertKostka}
Fix \(h\), and for \(1\leq j\leq s\), fix
\(\alpha^{(j)}\in\setN^h\) and \(\beta^{(j)}\in\setZ^h\).  Suppose that
\(N\alpha^{(j)}+\beta^{(j)}\) is nonnegative for \(N\gg0\), for each \(j\).
For \(N\gg0\), let \(w_j(N)\) be the permutation whose Lehmer code is obtained
from \(N\alpha^{(j)}+\beta^{(j)}\) by appending enough trailing zeroes to
satisfy the Lehmer-code inequalities.
For every fixed
\(\nu\in\setN^h\) and \(\eta\in\setZ^h\),
\[
[\xvec^{N\nu+\eta}]
\prod_{j=1}^s
\schubertS_{w_j(N)}(\xvec)
\]
is eventually polynomial in \(N\).  If the total degrees eventually disagree, the
coefficient is eventually zero.
\end{proposition}

\begin{proof}
For each factor, use the fixed orthodontic expression supplied by
Lemma~\ref{lem:stableOrthodonticData}.  The product assertion is then the last
part of Lemma~\ref{lem:fixedDemazureExpressions}, and the claim follows.
\end{proof}

We now use the Schubert analogue of the usual Weyl-denominator extraction of
Littlewood--Richardson coefficients.  For a weak composition
\(\alpha\in\setN^h\), write \(\defin{\schubertS_\alpha}\) for the Schubert polynomial
whose index has Lehmer code \(\alpha\).  Put
\(\defin{\mathbf 1}\coloneqq(1,\dotsc,1)\).  For \(\lambda\in\setZ^h\), choose an integer
\(k\) such that \(\lambda+k\mathbf 1\) is nonnegative, and define the
\defin{generalized Schubert polynomial}
\[
\defin{\schubertS_\lambda(\xvec)}
\coloneqq
(x_1\dotsm x_h)^{-k}\schubertS_{\lambda+k\mathbf 1}(\xvec).
\]
This is independent of \(k\).  Put
\[
\defin{\rho}\coloneqq(h-1,h-2,\dotsc,0),
\qquad
\defin{\Delta(\xvec)}\coloneqq\prod_{1\leq i<j\leq h}(x_i-x_j),
\]
and equip Laurent polynomials with the monomial pairing
\[
\defin{\langle\xvec^a,\xvec^b\rangle}
\coloneqq
\begin{cases}
1,& a=b,\\
0,& a\neq b.
\end{cases}
\]
For \(\lambda,\mu\in\setN^h\), Watanabe's duality formula
\cite[Lem.~7.4]{Watanabe2016} states that
\begin{equation}
\label{eq:watanabeDuality}
\left\langle
\schubertS_\lambda(\xvec),
\schubertS_{\rho-\mu}(\xvec^{-1})\Delta(\xvec)
\right\rangle
=
\begin{cases}
1,& \lambda=\mu,\\
0,& \lambda\neq\mu.
\end{cases}
\end{equation}

\begin{theorem}[Affine Schubert structure constants]
\label{thm:affineSchubertStructureConstants}
Fix \(\alpha,\gamma,\tau\in\setN^h\) and
\(\beta,\delta,\zeta\in\setZ^h\).  Suppose that
\(N\alpha+\beta\), \(N\gamma+\delta\), and \(N\tau+\zeta\) belong to
\(\setN^h\) for \(N\gg0\), and let \(u(N),v(N),w(N)\) be the permutations
with these respective Lehmer codes.  Then
\[
N\longmapsto c_{u(N),v(N)}^{w(N)}
\]
is eventually polynomial.
\end{theorem}

\begin{proof}
Unless the total degree of \(N\tau+\zeta\) equals the sum of the total degrees
of \(N\alpha+\beta\) and \(N\gamma+\delta\) identically in \(N\), the
coefficient is zero for all sufficiently large \(N\).  We therefore assume
that the degrees agree.
The Schubert polynomials indexed by weak compositions in \(\setN^h\) form a
basis of \(\setZ[x_1,\dotsc,x_h]\).  Pairing the Schubert expansion of
\(\schubertS_{u(N)}(\xvec)\schubertS_{v(N)}(\xvec)\) with
\eqref{eq:watanabeDuality} gives
\[
c_{u(N),v(N)}^{w(N)}
=
\left\langle
\schubertS_{u(N)}(\xvec)\schubertS_{v(N)}(\xvec),
\schubertS_{\rho-(N\tau+\zeta)}(\xvec^{-1})\Delta(\xvec)
\right\rangle.
\]
Since
\[
\Delta(\xvec)
=\sum_{\sigma\in\symS_h}
\operatorname{sgn}(\sigma)\xvec^{\sigma\rho},
\]
the monomial pairing turns this into
\begin{equation}
\label{eq:finiteSchubertExtraction}
c_{u(N),v(N)}^{w(N)}
=
\sum_{\sigma\in\symS_h}
\operatorname{sgn}(\sigma)
[\xvec^{\sigma\rho}]
\schubertS_{u(N)}(\xvec)
\schubertS_{v(N)}(\xvec)
\schubertS_{\rho-(N\tau+\zeta)}(\xvec).
\end{equation}

Choose an integer-valued affine function
\(\defin{q(N)}\coloneqq MN+K\) such that
\(\rho-(N\tau+\zeta)+q(N)\mathbf 1\) is nonnegative for \(N\gg0\).
It suffices to take \(M>\max_i\tau_i\) and then take \(K\) sufficiently
large.  For each fixed \(N\), apply the definition of generalized Schubert
polynomials with \(k=q(N)\).  Then each summand in
\eqref{eq:finiteSchubertExtraction} is
\[
[\xvec^{q(N)\mathbf 1+\sigma\rho}]
\schubertS_{u(N)}(\xvec)
\schubertS_{v(N)}(\xvec)
\schubertS_{\rho-(N\tau+\zeta)+q(N)\mathbf 1}(\xvec).
\]
The three Schubert indices and the extracted exponent are affine in \(N\).
Proposition~\ref{prop:productSchubertKostka} shows that every summand is
eventually polynomial.  There are only \(h!\) summands, independent of \(N\),
so their signed sum is eventually polynomial.  This completes the proof.
\end{proof}

\begin{corollary}[Stretched Schubert structure constants]
\label{cor:stretchedSchubertStructureConstants}
For fixed permutations \(u,v,w\), the function
\[
N\longmapsto
c_{\stretchperm{N}{u},\stretchperm{N}{v}}^{\stretchperm{N}{w}}
\]
is eventually polynomial.
\end{corollary}

\begin{proof}
Apply Theorem~\ref{thm:affineSchubertStructureConstants} with
\((\alpha,\beta)=(\code(u),0)\),
\((\gamma,\delta)=(\code(v),0)\), and
\((\tau,\zeta)=(\code(w),0)\).  This proves the corollary.
\end{proof}

This proves Conjecture~7.1 of I.~Pak and Z.~Slonim~\cite{PakSlonim2026x}.

\begin{corollary}[Stretched Littlewood--Richardson coefficients]
\label{cor:stretchedLittlewoodRichardson}
For fixed partitions \(\lambda,\mu,\nu\), the function
\[
N\longmapsto c_{N\lambda,N\mu}^{N\nu}
\]
is eventually polynomial.
\end{corollary}

\begin{proof}
Choose \(k\) such that \(\lambda,\mu,\nu\) have at most \(k\) parts, and
realize their Schur polynomials as the Grassmannian Schubert polynomials in
Lemma~\ref{lem:grassmannianStretching}.  After code stretching, these become
\(s_{N\lambda}\), \(s_{N\mu}\), and \(s_{N\nu}\), respectively.  Hence
\(c_{N\lambda,N\mu}^{N\nu}\) is a stretched Schubert structure constant, and
the claim follows from
Corollary~\ref{cor:stretchedSchubertStructureConstants}.
\end{proof}

The stronger assertion that \(c_{N\lambda,N\mu}^{N\nu}\) is given by a
polynomial for every \(N\geq0\) was proved by H.~Derksen and
J.~Weyman~\cite{DerksenWeyman2002}.  E.~Rassart later gave a chamber-complex
proof and established a piecewise-polynomial formula in the three partition
parameters~\cite{Rassart2004}.

\section*{Acknowledgments}

We thank Igor Pak for asking about polynomiality of stretched Schubert
coefficients; that question was the starting point for this paper.
OpenAI's Codex was used during the preparation of this paper for proof
exploration and editorial assistance.

\bibliographystyle{alpha}
\bibliography{key-schubert-kostka-ALEXANDERSSON-revised}

\end{document}
